\documentclass[10pt,a4paper]{amsart}
\usepackage[margin=19mm]{geometry}
\usepackage{amsmath,amssymb,mathtools}
\usepackage[scr=rsfs,cal=euler]{mathalfa}
\usepackage{xfrac}
\usepackage[unicode,hidelinks,bookmarks=false]{hyperref}
\newcommand{\ie}{i.e.\ }
\newcommand{\cf}{cf.\ }
\newcommand{\resp}{respectively\ }
\newcommand{\Subsection}{\subsection}
\providecommand{\nobreakdash}{\nobreak\hbox{-}}
\setlength{\emergencystretch}{2em}
\multlinegap0pt
\usepackage{amssymb}
\usepackage{xcolor}
\usepackage{cancel}
\usepackage{tikz}
\newtheorem{lemma}{Lemma}
\newcommand{\E}{\mathbb{E}}
\newcommand{\R}{\mathbb{R}}
\title{Nonlinear rough Fokker--Planck equations}
\author{Fabio Bugini, Peter K. Friz, Wilhelm Stannat}
\date{Source: arXiv:2507.17469v2 (CC BY 4.0)}
\begin{document}
\maketitle

\noindent\emph{Source and licence.} Nonlinear rough Fokker--Planck equations. Version arXiv:2507.17469v2 (CC BY 4.0), distributed by arXiv under the Creative Commons Attribution 4.0 International licence, http://creativecommons.org/licenses/by/4.0/, as recorded in the arXiv licence metadata for this exact version and shown on the abstract page https://arxiv.org/abs/2507.17469v2. Full licence notice: \texttt{licenses/arxiv-2507.17469-CC-BY-4.0.txt}.\par\smallskip

\noindent\textit{Editorial scope.} Counted unit: the article's lemma lemma:remainderinTaylorexpansion (source0.tex lines 333--339). The article's complete original proof of it is delivered with it (source0.tex lines 341--359, one \texttt{proof} environment of 19 lines). No mathematics is written, repaired or re-derived: the statement and the proof are the article's own lines, in the article's own notation, and no retained line is substituted or re-flowed.  No further source-local context is needed: every cross-reference of every retained line resolves inside the two delivered spans.  Nothing was omitted from any retained span, so the package carries no omission record.  No retained line contains a citation, so the wrapper carries no bibliography and no bibliography entry.  No retained line contains a figure environment, an image-inclusion command or drawing code, so no image is delivered and the package declares no asset dependency.  Editorial formatting only, in addition: an AMS \texttt{amsart} preamble that restates the article's own declarations and macros used by the retained lines, namely the environments \texttt{lemma} on separate counters, together with the standard packages those lines need.  The retained environments print in this document as Lemma 1 in order of appearance, whereas the article prints the same items with the numbering of its own full document.  The complete native source of the article as distributed in the arXiv e-print for this exact version is bundled unchanged as \texttt{sources/182082/source0.tex}.
    \begin{lemma} \label{lemma:remainderinTaylorexpansion}
        Let $g: \mathcal{P}_2(\R^d) \to \R$ be continuously L-differentiable and assume its L-derivative has a $\mu$-version $\partial_\mu g: \mathcal{P}_2(\R^d) \times \R^d \to \R^d$ which is Lipschitz continuous. Then, for any $\mu,\nu \in \mathcal{P}_2(\R^d)$ and for any coupling $\pi$ of $\mu$ and $\nu$, 
            \begin{equation*} 
            | \vartheta_{\mu,\nu} (\pi) | \le  |\partial_\mu g|_{Lip}  \int_{\R^d \times \R^d} |w-v|^2 \, \pi(dv, dw) .
            %| \vartheta_{\mu,\nu} (\pi) | \le 2 |\partial_\mu g|_{Lip} \mathcal{W}_2(\mu,\nu)^2 .
        \end{equation*}
    \end{lemma}

    \begin{proof}
        Let $\mu,\nu \in \mathcal{P}_2(\R^d)$ be fixed and let $X,Y$ be two arbitrary $\R^d$-valued random variables with $\mathcal{L}_X=\mu$ and $\mathcal{L}_Y = \nu$. 
        Being $\hat{g}:L_2(\Omega;\R^d) \to \R$ continuously Fréchet differentiable, by means of the mean value theorem we can write \begin{align*}
            g(\nu) - g(\mu) &= \hat{g}(Y) - \hat{g}(X) = \int_0^1 D\hat{g}(X + \theta(Y-X)) [Y-X] \, d\theta  \\ 
            &= \int_0^1 \E(\partial_\mu g(\mathcal{L}_{X + \theta(Y-X)}) (X + \theta(Y-X)) \cdot (Y-X) ) \, d\theta .
        \end{align*}
        Let $\pi$ be a coupling of $\mu$ and $\nu$. Then 
        \begin{equation*} 
            \begin{split}
                |\vartheta_{\mu,\nu} ( \pi)| &= |\hat{g}(Y) - \hat{g}(X) - \E(\partial_\mu g(\mu)(X) \cdot (Y-X))| \\
                &\le \int_0^1 \E(|(\partial_\mu g(\mathcal{L}_{X + \theta(Y-X)}) (X + \theta(Y-X)) - \partial_\mu g(\mu) (X) )\cdot (Y-X)|) \, d\theta \\
                &\le |\partial_\mu g|_{Lip} \int_0^1 \E((\mathcal{W}_2(\mathcal{L}_{X + \theta (Y-X)},\mu) + | \theta (Y-X)|)^2)^\frac12 \, d\theta \  \|Y-X\|_2 \\
                &\le |\partial_\mu g|_{Lip} \int_0^1 (2 \mathcal{W}_2(\mathcal{L}_{X + \theta (Y-X)},\mu)^2 + 2\theta^2 \E(|Y-X|^2))^\frac12 \, d\theta \ \|Y-X\|_2 \\
                &\le |\partial_\mu g|_{Lip} \|Y-X\|_2^2 =  |\partial_\mu g|_{Lip}  \int_{\R^d \times \R^d} |w-v|^2 \, \pi(dv,dw),
            \end{split}
        \end{equation*}
        where we used the fact that $\mathcal{W}_2(\mathcal{L}_{X + \theta (Y-X)},\mu) \le \| (X + \theta (Y-X)) - X\|_2 = \theta \|Y-X\|_2 $.
        %The conclusion follows from the arbitrariety of $X$ and $Y$.
    \end{proof}

\end{document}
